\documentclass[11pt]{article}
\usepackage{graphicx}
\usepackage{amsmath, amssymb, float}
\usepackage{mathtools}
\usepackage{enumitem}
\usepackage{url}
\usepackage{hyperref}
\usepackage{bm}
\allowdisplaybreaks

\DeclareMathOperator{\EX}{\mathbb{E}}% expected value
\DeclareMathOperator{\Var}{\operatorname{Var}}% variance
\DeclareMathOperator{\Span}{\operatorname{span}}% span
\DeclareMathOperator{\Ad}{\operatorname{Ad}}% Adjoint
\newcommand{\pd}[2]{\frac{\partial #1}{\partial #2}} % partial derivatives

\title{Spacecraft Attitude Control}
\author{Riley Bridges}

\begin{document}

\maketitle

\section{References}
\begin{enumerate}
  \item \href{https://space.skyrocket.de/doc_sdat/starlink-v2-mini.htm}{v2 mini satellite tracker}
  \item \href{https://starlink.com/technology}{starlink website (v2 mini component renders)}
  \item \href{https://control.asu.edu/Classes/MAE462/462Lecture15.pdf}{ASU ADCS lecture (high level starlink info)}
  \item \href{https://en.wikipedia.org/wiki/Starlink#Satellite_hardware}{starlink wikipedia}
  \item \href{https://en.wikipedia.org/wiki/List_of_Starlink_and_Starshield_launches}{List of Starlink launches}
  \item \href{https://www.mdpi.com/1424-8220/25/13/4079}{method for uncooperative tracking of starlink satellites}
  \item \href{https://apps.dtic.mil/sti/citations/AD0817503}{Desaturating reaction wheels using magnetorquers}
  \item \href{https://www.ietf.org/slides/slides-iepg-starlink-protocol-performance-00.pdf}{slides about starlink scheduling and performance}
  \item \href{https://digitalcommons.usu.edu/cgi/viewcontent.cgi?filename=0&article=4867&context=smallsat&type=additional}{paper about slew rates (maybe including starlink)}
  \item \href{https://www.teslarati.com/spacex-unveils-next-gen-starlink-v2-mini-satellites-ahead-of-monday-launch/}{v2 mini dimensions}
\end{enumerate}
\section{Mission Description}
We will be modeling and analyzing the Starlink V2-mini satellite. 
\subsection{Mission Objectives}
Thousands of these satellites are part of the starlink constellation, which provides global satellite internet and mobile broadband coverage.
\subsection{Orbit}
(Group 620 launched 8/8/2023) LEO 559km, $43^\circ$ inclination

\subsection{Sensors and Actuators}
The satellite is equipped with the following sensors:
\begin{itemize}
\item 3 axis gyroscope (assumed)
\item 3 axis accelerometer (assumed)
\item 3 axis magnetometer (assumed)
\item GNSS (assumed)
\item star tracker
\end{itemize}

and the following actuators:
\begin{itemize}
  \item Argon Hall effect thruster
  \item 4 reaction wheels
  \item magnetic torque rods
\end{itemize}

\subsection{ADCS Requirements}
0.1 degree pointing accuracy (totally made up), 6 deg/sec slew rate (possibly made up by google AI)

\section{Spacecraft Model}
The satellite's body frame is defined with its origin at the center of mass, the x-axis pointing in the direction the solar panels extend, the y-axis pointing along the perpendicular direction within the plane of the main body, and the z axis pointing along the direction that the antenna arrays are facing. This frame was chosen for simplicity while still aligning with the objective of pointing the antennas.

The satellite body was modeled as four rectangular prisms representing the main body, the two solar panels, and the beam connecting the solar panels to the main body. The dimensions of each component were drawn from Starlink filings and official renders. The mass of each component was estimated based on the total mass of the satellite and the relative volumes of each component.

For a uniform rectangular prism, the inertia tensor about its principle axes is defined as:
\begin{equation*}
  \mathcal{J} = \frac{m}{12}\begin{bmatrix}
    b^2+c^2 & 0 & 0 \\
    0 & a^2+c^2 & 0 \\
    0 & 0 & b^2 + c^2
  \end{bmatrix}
\end{equation*}

We can determine the satellite's total inertia tensor by modeling the body, solar panels, and beam as separate rectangular prisms, calculating their individual inertia tensors, and then applying the parallel axis theorem to shift them to a common center of mass before summing them together:
\begin{align*}
  \mathcal{J}_{total}^C &= \sum_i \mathcal{J}_i^C \\
  \mathcal{J}_i^C &= \mathcal{J}_i^B + m_i \left( r_i^\top r_i I - r_i r_i^\top \right)
\end{align*}
The satellite's center of mass can be calculated as the weighted average of the centers of mass of each component.

\section{Orbital Dynamics}
We can derive the satellite's orbital dynamics using Newton's law of gravitation:
\begin{align*}
  F_g &= \frac{GMm}{r^2} \\
  m\ddot{\mathbf{r}} &= -\frac{GMm}{\|\mathbf{r}\|^3} \mathbf{r} \\
  \mathbf{x} &= \begin{bmatrix} \mathbf{r} \\ \dot{\mathbf{r}} \end{bmatrix} \\
  \dot{\mathbf{x}} &= \begin{bmatrix} \dot{\mathbf{r}} \\ -\frac{GM}{\|\mathbf{r}\|^3} \mathbf{r} \end{bmatrix}
\end{align*}

We can then use our orbit parameters (satellite altitude $h$, earth radius $R_E$, and inclination $\iota$) to compute the satellite's initial state vector:
\begin{align*}
  \mathbf{r}_0 &= \begin{bmatrix} R_E + h \\ 0 \\ 0 \end{bmatrix} \\
  \mathbf{v}_0 &= \begin{bmatrix} 0 \\ \sqrt{\frac{GM}{R_E + h}} \cos\iota \\ \sqrt{\frac{GM}{R_E + h}} \sin\iota \end{bmatrix} \\
  \mathbf{x}_0 &= \begin{bmatrix} \mathbf{r}_0 \\ \mathbf{v}_0 \end{bmatrix}
\end{align*}
This orbit is not fully parameterized but this is enough for a rough recreation of what the satellite's orbit would look like.

\section{Attitude Dynamics}
Parameterizing the satellite's attitude with a quaternion, we can use Euler's equation and quaternion kinematics to derive the satellite's attitude dynamics:
\begin{align*}
  \dot{\mathbf{q}} &= \frac{1}{2} L(\mathbf{q}) H \boldsymbol{\omega}, \\ 
  L(\mathbf{q}) &= \begin{bmatrix}
    s & -\mathbf{v}^\top \\
    \mathbf{v} & sI + \mathbf{\hat{v}} \\
    \end{bmatrix}, \ 
    H = \begin{bmatrix}
      0 \\ I
  \end{bmatrix} \\
  \dot{\boldsymbol{\omega}} &= \mathcal{J}^{-1} \left( \boldsymbol{\tau} - \boldsymbol{\omega} \times \mathcal{J}\boldsymbol{\omega} \right) \\
\end{align*}
Where $\mathbf{q} = [s \ \mathbf{v}]^\top$ is the quaternion representing the satellite's orientation, $\boldsymbol{\omega}$ is the satellite's angular velocity expressed in the body frame, $\mathcal{J}$ is the inertia tensor expressed in the body frame, and $\boldsymbol{\tau}$ is the body frame torque vector.

The attitude dynamics were simulated using a 4th order Runge-Kutta integrator, with an initial angular velocity of 10 rpm about one principal axis plus a small initial perturbation. It was verified that the major and minor rotational axes were stable, while the intermediate axis was unstable.
\end{document}